\prefacesection{Abstract}
The minimization of deterministic finite automata (DFA) is a fundamental problem with practical applications in model checking, compilers, and verification. This paper implements and compares several massively parallel GPU algorithms for DFA minimization, including a transitive-closure (NC) method and three partitionrefinement variants (naive leader-election, sorting-based, and a hybrid partitionrefinement augmented with a partial transitive closure). All methods were implemented in CUDA and evaluated on synthetic benchmark families (Fibonacci and bit-splitter automata) as well as real-world models from the VLTS suite. Experimental results show that the transitive-closure approach, despite its theoretical appeal, requires prohibitive resources and does not scale well on current GPU hardware; by contrast, partition-refinement methods are more practical. The sorting-based partitioning achieves the best performance on many VLTS instances by enabling aggressive block splitting, while the hybrid enhancement produces large speedups on automata with long uniform paths. Implementation details, runtime, iteration counts, and memory usage are reported, and practical guidelines for selecting an appropriate GPU minimization strategy based on input structure and alphabet size are presented.
